\uuid{xrqm}
\exo7id{147}
\titre{exo7 147}
\auteur{ridde}
\organisation{exo7}
\datecreate{1999-11-01}
\isIndication{false}
\isCorrection{false}
\chapitre{Logique, ensemble, raisonnement}
\sousChapitre{Ensemble}
\cpge{mpsi-mp2i.raisonnement-vocabulaire-ensembliste.b, pcsi.raisonnement-vocabulaire-ensembliste.b, ptsi.raisonnement-vocabulaire-ensembliste.b, mpsi-mp2i.raisonnement-vocabulaire-ensembliste.c, pcsi.raisonnement-vocabulaire-ensembliste.d, ptsi.raisonnement-vocabulaire-ensembliste.d, bcpst1.vocabulaire-applications, bcpst1.logique-ensembles-raisonnement.b}
\module{Algèbre}
\niveau{L1}
\difficulte{}

\contenu{
\question{
Soit $E$ l'ensemble des fonctions de $\Nn$ dans $\left\{ 1, 2, 3\right\}$.
Pour $i = 1, 2, 3$ on pose $A_i = \left\{ f \in E / f (0) = i \right\}$.
Montrer que les $A_i$ forment une partition de $E$.
}
}
